\documentclass[12pt]{article}
\usepackage[margin=1in]{geometry}
\usepackage{mathptmx}
\usepackage[T1]{fontenc}
\usepackage[utf8]{inputenc}
\usepackage{setspace}
\usepackage{microtype}
\usepackage{fancyhdr}
\usepackage[hidelinks]{hyperref}
\pagestyle{fancy}
\fancyhf{}
\fancyfoot[C]{\thepage}
\renewcommand{\headrulewidth}{0pt}
\doublespacing
\setcounter{secnumdepth}{0}
\setlength{\parindent}{0.5in}
\setlength{\parskip}{0pt}
\setlength{\emergencystretch}{2em}
\begin{document}
\begin{center}
{\large\bfseries Carthage: Empire, Constitution, and the Sources}
\end{center}
\vspace{1em}
For nearly seven centuries Carthage was the greatest power in the
western Mediterranean. Yet no Carthaginian history survives. Everything
we know comes from Greek and Roman writers --- many of them its enemies
--- who worked from earlier sources now lost. Read critically, they
allow a coherent account of the city's rise, political order, and
destruction; read naively, they repeat the victors' prejudices.

\subsection{Origins and the commercial empire}

Appian reports two founding traditions: the founders were ``either Zorus
and Carchedon, or, as the Romans and the Carthaginians themselves think,
Dido, a Tyrian woman'' (\emph{Punic Wars} 1). The ox-hide story --- Dido
cuts the hide into a strip and encloses the hill that became the Byrsa,
``hide'' --- is an aetiology of the citadel, not a record of events.
Even the date was disputed: Appian puts the foundation fifty years
before the fall of Troy, while the traditional date is 814 BCE.

Carthaginian power is clearer. On the eve of the First Punic War,
Polybius says the Carthaginians ``had not only reduced Libya to
subjection, but a great part of Spain besides'' (\emph{Histories}
1.10.5). This was an empire of coasts, harbours, and tribute, not
annexed territory.

One text preserves a Carthaginian voice. The \emph{Periplus of Hanno} is
a Greek translation of the record that Hanno, king of the Carthaginians,
dedicated in the temple of Kronos. It records a state-sponsored
expedition beyond the Pillars of Heracles: ``It pleased the
Carthaginians that Hanno should voyage outside the Pillars of Hercules,
and found cities of the Libyphoenicians'' (\emph{Periplus} 1). The
narrative runs down the Atlantic coast of Africa, past a burning land,
to an island of \emph{Gorillai} (15--18).

\subsection{The constitution}

Carthage impressed Greek thinkers, whose analyses remain our best
evidence for its institutions. Aristotle counts Carthage among the three
well-ordered constitutions, with Sparta and Crete: ``The Carthaginians
are also considered to have an excellent form of government''
(\emph{Politics} 2.11). He notes kings, a council of elders, and a
popular assembly, and praises electing magistrates ``according to
merit.'' His strongest proof is stability: the Carthaginians ``have
never had any rebellion worth speaking of, and have never been under the
rule of a tyrant'' (2.11).

Carthage was, Polybius adds, ``originally well contrived,'' with kings,
an aristocratic council, and a sovereign people (\emph{Histories}
6.51.1--2). But by the Hannibalic War ``the Carthaginian constitution
had degenerated, and that of Rome was better'' (6.51.3): the multitude
had acquired ``the chief voice in deliberations,'' while at Rome the
senate retained it (6.51.6). His observation that Carthage relied on
``foreign and mercenary'' troops while Rome used citizens points to a
real structural difference (6.52.4).

\subsection{Sicily and the long contest with the Greeks}

Carthage's expansion repeatedly brought it into conflict with Greek
Sicily. When the Selinuntians pressed the Egestaeans, the latter
``dispatched ambassadors to Carthage, asking for aid and putting their
city in the hands of the Carthaginians,'' who were ``eager to acquire a
city so strategically situated'' (\emph{Library of History} 13.43.3--4).
After the sack of Selinus and Himera, he reports, they spared no
captives, crucifying some and inflicting ``unbearable outrages'' on
others (13.111.4). The struggle lasted two centuries: Dionysius attacked
the Carthaginian base at Motya (14.47.4), and in 310 BCE Agathocles
carried the war into Africa itself (20.3.1--3).

\subsection{Rome, Hannibal, and destruction}

Polybius makes the strategic logic of the First Punic War explicit: the
Romans feared that if the Carthaginians took Sicily they would be ``most
troublesome and dangerous neighbours, hemming them in on all sides''
(1.10.6). The war lasted twenty-four years and was ``the longest, most
unintermittent, and greatest war we know of'' (1.63.4--5). Carthage lost
Sicily.

Polybius records Hannibal's own account of the oath his father made him
swear at the altar, aged nine, ``never to be the friend of the Romans''
(3.11.7). Livy turns Hannibal into a moral portrait: a leader whose
``great merits were matched by great vices --- inhuman cruelty, a
perfidy worse than Punic'' (21.4). ``Punic faith'' was a Roman
stereotype for treachery, and Livy's characterization owes as much to
Roman moralizing as to his sources.

The end came in 146 BCE. Appian's account of the Third Punic War is our
fullest: the siege of Byrsa, the street fighting, the burning of the
city, and the suicide of Hasdrubal's wife (128--131). He records the
most famous scene: Scipio Aemilianus, watching Carthage burn, weeps and
quotes Homer on the fall of Troy, confessing his fear for Rome's own
future (132). Rome's relief was immense: ``no other war had so terrified
them at their own gates as the Punic wars'' (134).

\subsection{Conclusion}

The sources for Carthage are its enemies'. Aristotle uses the city to
test a theory of constitutions; Polybius to explain Rome's success; Livy
to define Roman virtue; Appian to satisfy a Greek audience fascinated by
Rome. Yet the texts agree that Carthage was a maritime commercial power,
governed by magistrates, a senate, and an assembly, reliant on
mercenaries and tribute, and long locked in conflict with Greek Sicily.
Reading the victors against the grain is the only way to recover the
city.

\newpage
\section*{Works Cited}

\noindent\hangindent=0.5in\hangafter=1 Appian of Alexandria. \emph{The Punic Wars} (\emph{Roman History}, Book
8.1). Translated by Horace White. Loeb Classical Library, 1912. Public
domain text at
\url{https://www.livius.org/sources/content/appian/appian-the-punic-wars/}.

\noindent\hangindent=0.5in\hangafter=1 Aristotle. \emph{Politics}. Translated by Benjamin Jowett. Internet
Classics Archive, MIT,
\url{https://classics.mit.edu/Aristotle/politics.2.two.html}.

\noindent\hangindent=0.5in\hangafter=1 Diodorus Siculus. \emph{Library of History}. Translated by C. H.
Oldfather and Russel M. Geer. Loeb Classical Library, 1933--1954. Public
domain text at LacusCurtius,
\url{https://penelope.uchicago.edu/Thayer/E/Roman/Texts/Diodorus_Siculus/home.html}.

\noindent\hangindent=0.5in\hangafter=1 Hanno. \emph{The Periplus of Hanno}. Translated by Wilfred H. Schoff.
Commercial Museum, Philadelphia, 1912. Public domain text at ToposText,
\url{https://topostext.org/work/247}.

\noindent\hangindent=0.5in\hangafter=1 Livy. \emph{The History of Rome}, Books 21--30. Translated by Rev.~Canon
Roberts. E. P. Dutton, 1912. Public domain text at ToposText,
\url{https://topostext.org/work/142}.

\noindent\hangindent=0.5in\hangafter=1 Polybius. \emph{The Histories}. Translated by W. R. Paton. Loeb
Classical Library, 1922. Public domain text at LacusCurtius,
\url{https://penelope.uchicago.edu/Thayer/E/Roman/Texts/Polybius/home.html}.
\end{document}
